% !TEX root = ../matan.tex

\section{Перестановки членов абсолютно и условно сходящихся рядов, теорема Римана}

\begin{Definition}{Перестановка ряда}{}
	\textbf{Перестановкой} натурального ряда называется биекция
	$\sigma:\NN\to\NN$. Ряд $\sum_{k=1}^{\infty} a_{\sigma(k)}$ называется
	\textbf{перестановкой} ряда $\sum_{k=1}^{\infty} a_k$ (тот же набор
	слагаемых, взятых в другом порядке).
\end{Definition}

\subsection*{Перестановка абсолютно сходящегося ряда}

\begin{Theorem}{Устойчивость абсолютно сходящегося ряда к перестановкам}{}
	Если ряд $\sum_{k=1}^{\infty} a_k$ сходится абсолютно к сумме $S$, то для
	любой перестановки $\sigma$ ряд $\sum_{k=1}^{\infty} a_{\sigma(k)}$ также
	сходится абсолютно, причём к той же сумме $S$.
\end{Theorem}

\begin{proof}
	Обозначим $S=\sum_{k=1}^{\infty} a_k$. Так как ряд $\sum|a_k|$ сходится, для
	любого $\eps>0$ найдётся $M$ такое, что хвост ряда модулей мал:
	\[
	\sum_{j=M+1}^{\infty}|a_j|<\eps.
	\]
	Так как $\sigma$ --- биекция, все номера $1,2,\dots,M$ встречаются среди
	$\sigma(1),\dots,\sigma(N)$ при достаточно большом $N$; выберем такое $N$.
	Тогда при любом $n\ge N$ в разности
	$\sum_{k=1}^{n} a_{\sigma(k)}-\sum_{j=1}^{M} a_j$ члены $a_1,\dots,a_M$
	сокращаются, и остаются лишь слагаемые $a_j$ с номерами $j>M$. Поэтому
	\[
	\left|\sum_{k=1}^{n} a_{\sigma(k)}-\sum_{j=1}^{M} a_j\right|
	\le\sum_{\substack{j>M}}|a_j|<\eps.
	\]
	С другой стороны, $\bigl|S-\sum_{j=1}^{M} a_j\bigr|\le\sum_{j>M}|a_j|<\eps$.
	По неравенству треугольника при $n\ge N$
	\[
	\left|\sum_{k=1}^{n} a_{\sigma(k)}-S\right|<2\eps,
	\]
	то есть $\sum a_{\sigma(k)}\to S$. Применяя уже доказанное к ряду $\sum|a_k|$
	(он абсолютно сходится к $\sum|a_k|$), получаем сходимость
	$\sum|a_{\sigma(k)}|$, то есть абсолютную сходимость переставленного ряда.
\end{proof}

\begin{Remark}{}{}
	Для условно сходящихся рядов это свойство теряется: перестановка может
	изменить сумму или вовсе разрушить сходимость. Именно поэтому абсолютную
	сходимость называют «безусловной». Точную меру произвола даёт следующая
	теорема.
\end{Remark}

\subsection*{Теорема Римана}

\begin{Theorem}{Теорема Римана о перестановках}{}
	Пусть ряд $\sum_{k=1}^{\infty} a_k$ сходится условно (то есть $\sum a_k$
	сходится, а $\sum|a_k|$ расходится). Тогда для любого числа $A\in\RR$
	существует перестановка $\sigma:\NN\to\NN$ такая, что
	$\sum_{k=1}^{\infty} a_{\sigma(k)}=A$. (Более того, можно добиться
	расходимости к $+\infty$ или к $-\infty$.)
\end{Theorem}

\begin{proof}
	Нулевые члены на сумму не влияют и могут быть вставлены позднее в любом
	месте, поэтому считаем $a_k\ne0$. Выпишем по порядку положительные и
	отрицательные члены исходного ряда:
	\[
	p_1,p_2,\dots\ (>0)\quad\text{--- положительные члены в их исходном порядке,}
	\]
	\[
	-q_1,-q_2,\dots\ (q_j>0)\quad\text{--- отрицательные члены в их исходном порядке.}
	\]

	\textbf{Шаг 1: $\sum p_j=+\infty$ и $\sum q_j=+\infty$.} Обозначим через
	$P_N$ сумму положительных, а через $-Q_N$ --- сумму отрицательных среди
	первых $N$ членов исходного ряда. Тогда частичная сумма и сумма модулей
	суть
	\[
	S_N=P_N-Q_N,\qquad \sum_{k=1}^{N}|a_k|=P_N+Q_N.
	\]
	Так как $\sum a_k$ сходится, $S_N=P_N-Q_N$ имеет конечный предел; так как
	$\sum|a_k|$ расходится, $P_N+Q_N\to+\infty$. Складывая и вычитая, получаем
	$2P_N=(P_N+Q_N)+S_N\to+\infty$ и $2Q_N=(P_N+Q_N)-S_N\to+\infty$. Значит,
	$\sum p_j=+\infty$ и $\sum q_j=+\infty$. Кроме того, из сходимости $\sum a_k$
	следует $a_k\to0$, поэтому $p_j\to0$ и $q_j\to0$.

	\textbf{Шаг 2: построение перестановки.} Зафиксируем $A\in\RR$. Строим
	переставленный ряд блоками. Сначала берём подряд положительные члены
	$p_1,p_2,\dots$ до тех пор, пока их сумма \emph{впервые} не превзойдёт $A$
	(это возможно, так как $\sum p_j=+\infty$). Затем добавляем подряд
	отрицательные члены $-q_1,-q_2,\dots$ до тех пор, пока текущая частичная
	сумма \emph{впервые} не станет меньше $A$ (возможно, так как
	$\sum q_j=+\infty$). Далее снова добавляем неиспользованные положительные
	члены до первого превышения $A$, затем отрицательные до первого опускания
	ниже $A$, и так далее. Так как каждый из запасов $\{p_j\}$, $\{q_j\}$
	неисчерпаем, процесс продолжается бесконечно и использует \emph{каждый}
	член ровно один раз; тем самым определена биекция $\sigma$.

	\textbf{Шаг 3: сходимость к $A$.} Обозначим частичные суммы построенного
	ряда через $\widetilde S_N$. Пусть на некотором шаге сумма впервые
	превысила $A$ добавлением положительного члена $p$; тогда до его добавления
	сумма была $\le A$, значит, превышение составляет $0\le\widetilde S-A\le p$,
	где $p$ --- этот последний добавленный член. Аналогично после
	отрицательного блока $0\le A-\widetilde S\le q$, где $q$ --- последний
	добавленный (по модулю) отрицательный член. Внутри блока частичные суммы
	монотонно движутся к точке перехода, поэтому отклонение $|\widetilde S_N-A|$
	в любой момент не превосходит модуля последнего члена, добавленного на
	текущем или предыдущем переходе.

	Поскольку $p_j\to0$ и $q_j\to0$, модули членов, на которых происходят
	переходы через уровень $A$, стремятся к нулю. Значит, для любого $\eps>0$
	найдётся $N_0$ такое, что при $N\ge N_0$ выполнено $|\widetilde S_N-A|<\eps$.
	Следовательно, $\widetilde S_N\to A$, то есть $\sum a_{\sigma(k)}=A$.
\end{proof}

\begin{Remark}{Смысл теоремы}{}
	У условно сходящегося ряда сумма зависит от порядка слагаемых, и притом
	максимально сильно: перестановкой можно получить любую вещественную сумму
	(а также $\pm\infty$ или вообще лишить ряд предела). Это прямо
	противоположно поведению абсолютно сходящихся рядов, сумма которых
	инвариантна относительно перестановок. Тем самым абсолютная сходимость ---
	это в точности «безусловная» сходимость.
\end{Remark}
